% Generated draft: requires mathematical and source review.
\documentclass[11pt]{article}
\usepackage[margin=25mm]{geometry}
\usepackage{fontspec}
\setmainfont{Latin Modern Roman}
\usepackage{amsmath,amssymb,mathtools,unicode-math}
\setmathfont{Latin Modern Math}
\usepackage{enumitem,needspace,xurl,hyperref}
\setlength{\parindent}{0pt}
\setlength{\parskip}{5pt}
\setlength{\emergencystretch}{3em}
\newcommand{\N}{\mathbb{N}}
\newcommand{\Z}{\mathbb{Z}}
\newcommand{\Q}{\mathbb{Q}}
\newcommand{\R}{\mathbb{R}}
\newcommand{\C}{\mathbb{C}}
\newcommand{\F}{\mathbb{F}}
\newcommand{\A}{\mathbb{A}}
\newcommand{\PP}{\mathbb{P}}
\newcommand{\E}{\mathbb{E}}
\newcommand{\eps}{\varepsilon}
\newcommand{\dd}{\,\mathrm{d}}
\begin{document}
\section*{UnsolvedMath ID: 2803}
\textbf{Source problem code:} catalogue entry
\subsection{Definitions and mathematical statement}
A closed hyperbolic $3$-manifold is a compact boundaryless manifold with a complete metric of sectional curvature $-1$. Its global injectivity radius is the infimum of its pointwise injectivity radii, equivalently half its shortest nontrivial closed-geodesic length. In the usual orientable convention, a closed irreducible $3$-manifold is Haken if it contains a closed embedded two-sided incompressible surface of positive genus; non-Haken means no such surface. The source does not explicitly fix orientability or all foliation-regularity conventions.

A cofinal tower of regular finite-sheeted connected covers is
\[
\cdots\longrightarrow M_2\longrightarrow M_1\longrightarrow M_0=M,
\]
where the associated nested finite-index subgroups $G_i=\pi_1(M_i)\leq\pi_1(M)$ are normal in $\pi_1(M)$ and $\bigcap_iG_i=\{1\}$. The source requires regularity over the base, not only at adjacent stages.

The three displayed parts are:
\begin{enumerate}
\item For every real $R>0$, does there exist a closed non-Haken hyperbolic $3$-manifold $M$ with $\operatorname{inj}(M)>R$?
\item Does there exist a cofinal regular tower over a closed hyperbolic $3$-manifold in which every manifold is non-Haken?
\item Are there towers of hyperbolic rational homology $3$-spheres satisfying each of the following proposed alternatives: all fundamental groups are not left-orderable; all manifolds admit no coorientable taut foliation; or all manifolds are L-spaces? The source does not state whether one common tower is required for all three, nor whether ``tower'' here inherits cofinality and regularity from part (b). Those choices must be explicit; otherwise constant towers or identity covers are not excluded by the literal word alone.
\end{enumerate}
A rational homology sphere has the same rational homology as $S^3$. A left order is a total group order preserved by left multiplication. A coorientable codimension-one foliation has a globally oriented normal line bundle; tautness means that every leaf meets a closed transversal in the usual closed-manifold setting. An L-space, in the standard $\F_2$ coefficient convention, is a rational homology sphere $M$ with
\[
\dim_{\F_2}\widehat{HF}(M;\F_2)=|H_1(M;\Z)|,
\]
where $\widehat{HF}$ is hat Heegaard Floer homology. The relation among the three properties is the separately cited L-space conjecture, not an assumed theorem making them equivalent here.

The background attributes part (a) to Cooper and Kirby's 1997 Problem 3.58. It reports that injectivity radii tend to infinity in a cofinal regular tower over a closed hyperbolic manifold, so a negative answer to (a) implies a negative answer to (b). Since such towers exist, a negative answer to either would give an alternative route to virtual Hakenness; this is motivation, not a replacement for the stronger non-Haken questions, and the Virtual Haken Conjecture is already reported proved by Agol.

The source reports Frank Calegari--Dunfield sequences of closed hyperbolic rational homology spheres with injectivity radii tending to infinity, including a tower covering the Weeks manifold as a candidate for part (b). Rational homology-sphere status does not establish non-Hakenness. It explicitly says these examples are not integral homology spheres and records an expectation, referred to Problem 3.6, against cofinal regular towers entirely of hyperbolic integral homology spheres. That expectation is not asserted as proved. It also reports Danny Calegari--Dunfield's non-left-orderability theorem for the Weeks group and separately asks whether the groups in Frank Calegari--Dunfield's tower share that property; non-left-orderability is not silently inherited by subgroups.

The registry is open and the supplied triage dated 2026-08-17 reports partial progress, retaining the non-Haken and tower alternatives unresolved. Its shorthand reference to finite-quotient questions does not supply an additional displayed mathematical subproblem. All theorem, example, and status claims above are supplied reports rather than independent verification here.
\subsection{Short English statement}
Can non-Haken closed hyperbolic $3$-manifolds have arbitrarily large injectivity radius, and can every level of a cofinal regular tower remain non-Haken? A further part asks for rational-homology-sphere towers with orderability, foliation, or L-space properties. Known large-radius rational homology spheres do not settle those extra conditions.
\subsection{Sources}
\begin{itemize}
\item K3: A New Problem List in Low-Dimensional Topology, Problem 3.5 (2026), supplied Kirby-style source: \url{https://aimath.org/pastworkshops/kirbylistrep.pdf}. This citation is retained without opening or generating a PDF.
\item [Kir97] R. C. Kirby, ``Problems in low-dimensional topology,'' in W. Kazez, editor, Geometric Topology, American Mathematical Society/International Press, Providence (1997), original Problem 3.58 for Cooper's question.
\item [Ago13] Ian Agol, ``The virtual Haken conjecture,'' Documenta Mathematica 18 (2013), 1045--1087, with appendix by Agol, Daniel Groves, and Jason Manning, DOI 10.4171/DM/421: \url{https://elibm.org/article/10000267}.
\item [CD06] Frank Calegari and Nathan M. Dunfield, ``Automorphic forms and rational homology 3-spheres,'' Geometry and Topology 10 (2006), 295--329, DOI 10.2140/gt.2006.10.295.
\item [CD03] Danny Calegari and Nathan M. Dunfield, ``Laminations and groups of homeomorphisms of the circle,'' Inventiones Mathematicae 152(1) (2003), 149--204, DOI 10.1007/s00222-002-0271-6.
\item Local registry data/unsolvedmath/problems.json, permanent catalogue entry, source code catalogue entry, ``Kirby Problem 3.5''; all three parts and six literature notes, and literature triage dated 2026-08-17. The source also refers to Problems 3.6 and 3.48. Dataset attribution: Ulam AI and the UnsolvedMath Contributors.
\end{itemize}
\end{document}
